\documentclass[11pt,a4paper]{article}
\usepackage[margin=25mm]{geometry}
\usepackage{amsmath,amssymb,bm,booktabs,graphicx,hyperref}
\hypersetup{colorlinks=true,linkcolor=blue,urlcolor=blue}
\newcommand{\R}{\mathbb{R}}
\newcommand{\SO}{\mathrm{SO}(3)}
\newcommand{\hatm}[1]{[#1]_\times}
\newcommand{\T}{^{\!\top}}
\newcommand{\e}{\bm{e}}
\title{Model, estimation and control of the air--ground mobile cable-driven parallel robot}
\author{Derivations for the C++ core (\texttt{core/})}
\date{}
\begin{document}
\maketitle

\noindent This document derives the model that \texttt{core/} implements, from the single unit (one cable, one
winch, one quadrotor, one ground robot) to the coupled machine, and then the estimators and controllers built on
it. Each section names the function that implements it. Numbers quoted in Section~\ref{sec:results} come from
\texttt{tests/core\_test.py} and \texttt{tests/isaac\_physics\_test.py}.

\tableofcontents

\section{Frames and notation}
World frame $\mathcal{W}$: $z$ up, ground at $z=0$, $\e_3=(0,0,1)\T$. Every rigid body carries a frame at its
centre of mass; $\bm{R}\in\SO$ maps body to world. $\bm{\Omega}$ is the angular velocity in the body frame, so
$\dot{\bm{R}}=\bm{R}\hatm{\bm{\Omega}}$, with $\hatm{\bm a}\bm b=\bm a\times\bm b$. Indices: cables
$i=1..8$; cables $1..4$ go from the platform's top corners to the quadrotors, $5..8$ from the bottom corners to the
ground robots.

\begin{center}\resizebox{\ifdim\width>\textwidth\textwidth\else\width\fi}{!}{\begin{tabular}{ll}
\toprule
$\bm p,\bm v,\bm R,\bm\Omega$ & platform position, velocity, attitude, body rate \\
$m,\bm J$ & platform mass and inertia (3 kg, cube of side 0.30 m) \\
$\bm b_i$ & attachment point of cable $i$ in the platform frame (cube corners) \\
$\bm a_i,\dot{\bm a}_i$ & anchor (fairlead) position and velocity of cable $i$ \\
$d_i,\bm u_i$ & chord length and unit vector from attachment point to anchor \\
$L_i,T_i$ & unstretched cable length, cable tension \\
$EA,c_{EA},\rho$ & cable axial stiffness, damping, linear density \\
\bottomrule
\end{tabular}}\end{center}

\section{Unit models}
\subsection{Cable}\label{sec:cable}
\paragraph{Kinematics} (\texttt{cable\_geometry}, \texttt{cable\_rates}). With $\bm r_i=\bm R\bm b_i$,
\begin{equation}
 \bm d_i=\bm a_i-\bm p-\bm r_i,\qquad d_i=\|\bm d_i\|,\qquad \bm u_i=\bm d_i/d_i,\qquad
 \dot d_i=\bm u_i\T\!\left(\dot{\bm a}_i-\bm v-\bm\omega\times\bm r_i\right),\quad \bm\omega=\bm R\bm\Omega .
 \label{eq:kin}
\end{equation}
\paragraph{Elasticity} (\texttt{cable\_tension}). The cable is a tension-only Kelvin--Voigt element whose
stiffness and damping scale with $1/L_i$:
\begin{equation}
 T_i=\begin{cases}\max\!\left(0,\ \dfrac{EA}{L_i}(d_i-L_i)+\dfrac{c_{EA}}{L_i}(\dot d_i-\dot L_i)\right) & d_i>L_i\\[2mm]
 0 & d_i\le L_i \end{cases}
 \label{eq:kv}
\end{equation}
Statically $d_i=L_i(1+T_i/EA)$, which is how a winch encoder ($L_i$) and a load cell ($T_i$) together measure the
chord (\texttt{chord\_length}).
\paragraph{Sag.} A cable of weight per length $\rho g$ at elevation $\beta_i$ ($\sin\beta_i=u_{i,z}$) carries a
normal load $w_n=\rho g\cos\beta_i$. For small sag the profile is a parabola and the chord is shorter than the
stretched length $s_i=L_i(1+T_i/EA)$ by
\begin{equation} d_i=s_i\left(1-\frac{(w_n s_i)^2}{24\,T_i^2}\right). \end{equation}
For the 2~mm rope ($\rho=2.6$ g/m) at 5~m and 15~N this is 0.04~mm, so the simulation uses straight cables with
half of each cable's weight lumped at its ends; the correction is available in the observer
(\texttt{observer.sag}).

\subsection{Winch}\label{sec:winch}
Drum of radius $r_d$, reflected inertia $J_w$, viscous and Coulomb friction, in cable coordinates
($m_w=J_w/r_d^2=5$ kg) (\texttt{winch\_step}):
\begin{equation}
 m_w\ddot L_i=F_i+T_i-c_v\dot L_i-F_c\,\mathrm{sgn}(\dot L_i),\qquad |F_i|\le F_{\max}\Bigl(1-\tfrac{|\dot L_i|}{v_{\max}}\Bigr)\ \text{when driving.}
\end{equation}
The motor force is one of two servos, both with the tension feed-forward $T^{\mathrm{ff}}_i$:
\begin{align}
 \text{length mode:}\quad F_i&=-T^{\mathrm{ff}}_i+m_w\!\left[\omega_w^2(L^c_i-L_i)+2\zeta_w\omega_w(\dot L^c_i-\dot L_i)\right], \label{eq:lenmode}\\
 \text{tension mode:}\quad F_i&=-T^{\mathrm{ff}}_i+c_w(\dot L^c_i-\dot L_i),\qquad c_w=2\zeta_w\omega_w m_w=450\ \mathrm{N\,s/m}. \label{eq:tenmode}
\end{align}
In tension mode at steady speed $T_i=T^{\mathrm{ff}}_i+c_w(\dot L_i-\dot L^c_i)+c_v\dot L_i+F_c\,\mathrm{sgn}\dot L_i$:
the winch is a force source in series with a damper referenced to the commanded rate.

\subsection{Quadrotor}\label{sec:quad}
Mass $m_q$, inertia $\bm J_q$, four rotors with thrust $f_j=k_f\varpi_j^2$, first-order speed dynamics
$\dot\varpi_j=(\varpi^c_j-\varpi_j)/\tau_m$, and mixer (\texttt{quad\_mixer})
\begin{equation}
\begin{pmatrix}f\\ \bm M\end{pmatrix}=\begin{pmatrix}1&1&1&1\\ y_1&y_2&y_3&y_4\\ -x_1&-x_2&-x_3&-x_4\\ -\kappa s_1&-\kappa s_2&-\kappa s_3&-\kappa s_4\end{pmatrix}\!\begin{pmatrix}f_1\\f_2\\f_3\\f_4\end{pmatrix}.
\end{equation}
The cable pulls the fairlead (offset $\bm r_a$ from the centre of mass) with $-T\bm u$:
\begin{align}
 m_q\dot{\bm v}_q&=f\bm R_q\e_3-m_qg\e_3-T\bm u-\tfrac12\rho gL\,\e_3+\bm F_{\mathrm{drag}}, \label{eq:quadtrans}\\
 \bm J_q\dot{\bm\Omega}_q&=\bm M-\bm\Omega_q\times\bm J_q\bm\Omega_q+\bm r_a\times\bm R_q\T(-T\bm u), \label{eq:quadrot}
\end{align}
with $\bm F_{\mathrm{drag}}=-\bm R_q\,\mathrm{diag}(k_d,k_d,0)\bm R_q\T\bm v_q-c_d\|\bm v_q\|\bm v_q$ (\texttt{quad\_wrench}).

\subsection{Ground robot}\label{sec:ugv}
Differential drive, pose $(x,y,\theta)$, forward speed $v$, yaw rate $\omega$, wheel radius $r$, track $b$.
Rolling without slipping gives $\dot x=v\cos\theta$, $\dot y=v\sin\theta$, $\dot\theta=\omega$ and wheel speeds
$\varpi_{l,r}=(v\mp\tfrac b2\omega)/r$. The wheels are velocity servos, $\tau_{l,r}=k_w(\varpi^c_{l,r}-\varpi_{l,r})$,
$|\tau|\le\tau_{\max}$. The fairlead sits on the axle axis at height $h$, so the cable force
$\bm F=-T\bm u$ has no yaw moment:
\begin{equation}
 \Bigl(m_g+\tfrac{2J_{wh}}{r^2}\Bigr)\dot v=\frac{\tau_l+\tau_r}{r}+\bm F\!\cdot\!\bm h,\qquad
 \Bigl(I_z+\tfrac{J_{wh}b^2}{2r^2}\Bigr)\dot\omega=\frac{(\tau_r-\tau_l)\,b}{2r},
\end{equation}
with $\bm h=(\cos\theta,\sin\theta,0)\T$. The lateral component $\bm F\cdot\bm h_\perp$ and the drive force are
carried by tyre friction; the model reports the friction demand over $\mu N$, $N=m_gg+F_z$ (\texttt{Simulator::traction\_use}).

\subsection{Platform}\label{sec:platform}
Newton--Euler about the centre of mass with the external wrench $(\bm f_e,\bm\tau_e)$ of the tool and payload:
\begin{align}
 m\dot{\bm v}&=\sum_i T_i\bm u_i-mg\e_3+\bm f_e, &
 \bm J\dot{\bm\Omega}&=\bm R\T\Bigl(\sum_i\bm r_i\times T_i\bm u_i+\bm\tau_e\Bigr)-\bm\Omega\times\bm J\bm\Omega .
 \label{eq:plat}
\end{align}

\section{The coupled machine}
\subsection{Structure matrix}
Collecting \eqref{eq:plat}, the cables apply the wrench $\bm W\bm t$ with
\begin{equation}
 \bm W(\bm p,\bm R,\bm a)=\begin{pmatrix}\bm u_1&\cdots&\bm u_8\\ \bm r_1\times\bm u_1&\cdots&\bm r_8\times\bm u_8\end{pmatrix}\in\R^{6\times 8},\qquad \bm t\ge 0 .
\end{equation}
By \eqref{eq:kin}, $\dot{\bm d}=-\bm W\T\bm\xi+\bm U\T\dot{\bm a}$ with the platform twist $\bm\xi=(\bm v,\bm\omega)$:
$-\bm W\T$ is the inverse-kinematics Jacobian. Eight cables on six degrees of freedom leave a two-dimensional
null space for internal tension.

\subsection{State and dimension}
The plant state (\texttt{Simulator}) is
\[
 \underbrace{(\bm p,\bm v,\bm R,\bm\Omega)}_{12}\ \oplus\ \underbrace{4\times(\bm p_q,\bm v_q,\bm R_q,\bm\Omega_q,\bm\varpi)}_{4\times 16}\ \oplus\ \underbrace{4\times(x,y,\theta,v,\omega)}_{4\times 5}\ \oplus\ \underbrace{(\bm L,\dot{\bm L})}_{16}
\]
i.e.\ 112 states, coupled only through the eight cable tensions \eqref{eq:kv}: each tension depends on the
platform pose, one robot's pose and one winch. The inputs are 16 rotor speed commands, 8 wheel speed commands and
8 winch commands $(L^c,\dot L^c,T^{\mathrm{ff}},\text{mode})$. Integration is semi-implicit Euler at 1~kHz, the
same step as PhysX.

\subsection{Tension bounds from the robots' limits}\label{sec:caps}
Each robot holds exactly one cable, so its limits are bounds on that one tension.

\paragraph{Quadrotor} (\texttt{drone\_tension\_cap}). In equilibrium \eqref{eq:quadtrans} needs thrust
$\bm f=m_qg\e_3+T\bm u$. Tilt $\le\theta_{\max}$ and $\|\bm f\|\le f_{\max}$ give
\begin{equation}
 T\le\frac{m_qg\tan\theta_{\max}}{u_{xy}-u_z\tan\theta_{\max}}\ \ (\text{if }u_{xy}>u_z\tan\theta_{\max}),\qquad
 T\le-m_qg\,u_z+\sqrt{(m_qg\,u_z)^2-(m_qg)^2+f_{\max}^2},
 \label{eq:capdrone}
\end{equation}
with $u_{xy}=\|(u_x,u_y)\|$.

\paragraph{Ground robot} (\texttt{ugv\_tension\_cap}). The cable pulls with $T u_{xy}$ horizontally and
$-Tu_z$ vertically. $\bm u$ points from the platform to the anchor, so for a ground robot $u_z<0$ and the cable
lifts it: the normal force is $N=m_gg+Tu_z$. Sliding: $Tu_{xy}\le\mu(m_gg+Tu_z)$. Tipping about the edge of the support
polygon at distance $a$ (inradius of the wheel--caster rhombus, 0.148 m) with the fairlead at height $h$:
$Tu_{xy}h\le(m_gg+Tu_z)a$. Hence
\begin{equation}
 T\le\frac{\mu m_gg}{u_{xy}-\mu u_z},\qquad T\le\frac{m_gg\,a}{u_{xy}h-u_za}.
 \label{eq:capugv}
\end{equation}
For the default layout the tipping bound is 63~N at 80\% margin, just above the 60~N winch limit; the sliding
bound is about 130~N.

\subsection{Tension distribution}\label{sec:td}
Given the wrench $\bm w$ the cables must supply, the tensions solve (\texttt{tension\_distribution})
\begin{equation}
 \min_{\bm t}\ \lambda\|\bm t-t_{\mathrm{ref}}\bm 1\|^2+\lambda\kappa\|\bm t-\bm t^-\|^2+\|\bm W\bm t-\bm w\|^2
 \quad\text{s.t.}\quad \bm t_{\min}\le\bm t\le\bm t_{\max}(\bm u),
 \label{eq:qp}
\end{equation}
where $\bm t_{\max}$ is the smaller of the winch limit and \eqref{eq:capdrone}, \eqref{eq:capugv}. It is a strictly
convex box-constrained QP in eight variables, solved exactly by a primal active-set method
(\texttt{solve\_box\_qp}): minimise over the free set, step to the first bound hit, release the bound whose
multiplier has the wrong sign, stop when the projected gradient vanishes. Without active bounds and $\kappa=0$ the
solution is the damped pseudo-inverse $\bm t=\bm t_{\mathrm{ref}}+\bm W\T(\bm W\bm W\T+\lambda\bm I)^{-1}(\bm w-\bm W\bm t_{\mathrm{ref}})$.
The damping follows the conditioning of $\bm W$ (\texttt{adaptive\_lambda}):
\begin{equation}
 \lambda=\lambda_{\min}+\bigl(1-(\sigma_{\min}/\sigma_\epsilon)^2\bigr)(\lambda_{\max}-\lambda_{\min})\ \ \text{for }\sigma_{\min}(\bm W)<\sigma_\epsilon,\qquad\lambda=\lambda_{\min}\ \text{otherwise,}
\end{equation}
so the wrench is met to $10^{-6}$ in normal layouts ($\sigma_{\min}=0.27$ for the crossed layout) and is not
forced through vanishing moment arms near a singular one.

\section{Sensors and estimation}\label{sec:est}
\begin{center}\resizebox{\ifdim\width>\textwidth\textwidth\else\width\fi}{!}{\begin{tabular}{llll}
\toprule
Sensor & Measures & Rate & Noise (1$\sigma$) \\
\midrule
Drone IMU & specific force, body rate & 500 Hz & $6.9\!\times\!10^{-4}$ m/s$^2$/$\sqrt{\text{Hz}}$, $4.9\!\times\!10^{-5}$ rad/s/$\sqrt{\text{Hz}}$, biases \\
Drone fix & position, heading & 50 Hz & 5 mm, 0.5$^\circ$ \\
Wheel encoders & wheel rates & 100 Hz & 0.02 rad/s \\
Ground-robot fix (lidar) & $x,y,\theta$ & 15 Hz & 5 mm, 0.3$^\circ$ \\
Winch encoder & $L_i,\dot L_i$ & 1 kHz & 0.03 mm, 1 mm/s \\
Winch load cell & $T_i$ & 1 kHz & 0.1 N \\
Robot camera + AprilTag & platform pose in the camera frame & 30 Hz & 3 mm, 0.3$^\circ$ \\
\bottomrule
\end{tabular}}\end{center}
The measurements are generated from the true simulator state by \texttt{SensorSuite} (rates, white noise, IMU
bias and bias random walk). The camera images and lidar scans themselves are not processed in this loop.

\subsection{Quadrotor: error-state Kalman filter}
Nominal state $(\bm p_q,\bm v_q,\bm R_q,\bm b_g,\bm b_a)$, propagated with the IMU sample $(\bm a_m,\bm\omega_m)$:
$\dot{\bm p}_q=\bm v_q$, $\dot{\bm v}_q=\bm R_q(\bm a_m-\bm b_a)-g\e_3$, $\dot{\bm R}_q=\bm R_q\hatm{\bm\omega_m-\bm b_g}$. The error
$\delta\bm x=(\delta\bm p,\delta\bm v,\delta\bm\theta,\delta\bm b_g,\delta\bm b_a)$, with $\bm R_q\leftarrow\bm R_q\exp\hatm{\delta\bm\theta}$, obeys
\begin{equation}
 \delta\dot{\bm p}=\delta\bm v,\quad \delta\dot{\bm v}=-\bm R_q\hatm{\bm a_m-\bm b_a}\delta\bm\theta-\bm R_q\delta\bm b_a,\quad
 \delta\dot{\bm\theta}=-\hatm{\bm\omega_m-\bm b_g}\delta\bm\theta-\delta\bm b_g .
\end{equation}
The position fix updates with $\bm H=[\bm I\ \bm 0\ \bm 0\ \bm 0\ \bm 0]$; the heading
$\psi=\mathrm{atan2}(R_{21},R_{11})$ with $\partial\psi/\partial\delta\bm\theta=(R_{11}\bm A_2-R_{21}\bm A_1)/(R_{11}^2+R_{21}^2)$,
$\bm A=-\bm R_q\hatm{\e_1}$ (\texttt{DroneEskf}). The fairlead is $\bm a=\bm p_q+\bm R_q\bm r_a$.

\subsection{Ground robot: odometry + pose fix}
$(x,y,\theta)$ is propagated with the encoder odometry $v=\tfrac r2(\varpi_l+\varpi_r)$, $\omega=\tfrac rb(\varpi_r-\varpi_l)$ and
corrected by the pose fix (\texttt{UgvEkf}); the anchor is $(x,y,h)$.

\subsection{Platform: tension-driven observer}\label{sec:obs}
The platform carries no inertial sensor. Its observer (\texttt{PlatformObserver}) uses the model
\eqref{eq:plat} itself as the predictor, driven by the \emph{measured} tensions, and adds the unknown external
wrench as a random-walk state: nominal state $(\bm p,\bm v,\bm R,\bm\Omega,\bm d_f,\bm d_\tau)$,
\begin{align}
 \dot{\bm v}&=\tfrac1m\Bigl(\sum_iT_i\bm u_i+\bm f_k+\bm d_f\Bigr)-g\e_3, &
 \bm J\dot{\bm\Omega}&=\sum_i\bm b_i\times\bm R\T T_i\bm u_i+\bm R\T\bm\tau_k+\bm d_\tau-\bm\Omega\times\bm J\bm\Omega,
\end{align}
with $(\bm f_k,\bm\tau_k)$ the known tool wrench (wrist load cell). Because $T_i$ is measured, the stiff elastic
law \eqref{eq:kv} does not appear in the predictor; only the direction $\bm u_i$ depends on the state:
$\partial(T_i\bm u_i)/\partial\bm x_{a,i}=-\tfrac{T_i}{d_i}\bm P_i$, $\bm P_i=\bm I-\bm u_i\bm u_i\T$, the geometric stiffness. With
$\delta\bm x_{a,i}=\delta\bm p-\bm R\hatm{\bm b_i}\delta\bm\theta$ the error dynamics are
\begin{align}
 m\,\delta\dot{\bm v}&=-\sum_i\tfrac{T_i}{d_i}\bm P_i\,\delta\bm p+\sum_i\tfrac{T_i}{d_i}\bm P_i\bm R\hatm{\bm b_i}\,\delta\bm\theta+\delta\bm d_f,\\
 \bm J\,\delta\dot{\bm\Omega}&=-\sum_i\hatm{\bm b_i}\bm R\T\tfrac{T_i}{d_i}\bm P_i\,\delta\bm p
 +\sum_i\Bigl(T_i\hatm{\bm b_i}\hatm{\bm R\T\bm u_i}+\hatm{\bm b_i}\bm R\T\tfrac{T_i}{d_i}\bm P_i\bm R\hatm{\bm b_i}\Bigr)\delta\bm\theta\notag\\
 &\quad+\bigl(\hatm{\bm J\bm\Omega}-\hatm{\bm\Omega}\bm J\bigr)\delta\bm\Omega+\delta\bm d_\tau,
\end{align}
plus $\delta\dot{\bm p}=\delta\bm v$, $\delta\dot{\bm\theta}=-\hatm{\bm\Omega}\delta\bm\theta+\delta\bm\Omega$.

\paragraph{Cable-length update (tension-based localisation).} Each taut cable ($T_i>1$ N) gives the scalar
measurement $z_i=L_i(1+T_i/EA)=\|\bm a_i-\bm p-\bm R\bm b_i\|$ with
\begin{equation}
 \bm H_i=\bigl[-\bm u_i\T\ \ \bm 0\ \ \bm u_i\T\bm R\hatm{\bm b_i}\ \ \bm 0\ \ \bm 0\ \ \bm 0\bigr],
\end{equation}
i.e.\ row $i$ of $-\bm W\T$ in body coordinates: the eight lengths observe the full pose whenever $\bm W$ has rank
6. The anchors come from the robots' own filters, so the measurement variance includes the anchor error.
\paragraph{Tag update.} A camera on robot $k$ measures the platform pose in its own frame,
$(\bm p_{\mathrm{rel}},\bm R_{\mathrm{rel}})$; with the robot's estimated pose
$\bm p_{\mathrm{tag}}=\hat{\bm p}_k+\hat{\bm R}_k\bm p_{\mathrm{rel}}$, $\bm R_{\mathrm{tag}}=\hat{\bm R}_k\bm R_{\mathrm{rel}}$, and the residuals are
$\bm p_{\mathrm{tag}}-\bm p$ and $\log(\bm R\T\bm R_{\mathrm{tag}})^\vee$. An attitude error $\epsilon$ of a robot at range $\ell$
moves the tag position by $\ell\epsilon$ (7 mm for 0.1$^\circ$ at 4 m), so the tags are weighted well below the
cable lengths; they guard against a wrong mode rather than set the precision.
\paragraph{Why not length control alone.} A length-controlled cable turns an anchor position error $\delta a$
along the cable into a tension error
\begin{equation} \delta T=\frac{EA}{L}\,\delta a\qquad(34\ \text{N/mm at }L=2.9\ \text{m}), \end{equation}
so anchors known to a few millimetres cannot be combined with eight stiff length servos. With four tension-mode
cables the four length-mode cables only fix four of the six platform degrees of freedom, an anchor error moves
the platform instead of loading the cables, and the tensions stay near their set-points. This is the estimation
side of the hybrid scheme of Section~\ref{sec:hybrid}.

\section{Control}
\subsection{Platform tracking on SE(3)}
Reference $(\bm p_d,\bm v_d,\bm a_d,\bm R_d,\bm\Omega_d)$. Errors
\begin{equation}
 \e_p=\bm p_d-\bm p,\quad \e_v=\bm v_d-\bm v,\quad \e_R=\tfrac12(\bm R_d\T\bm R-\bm R\T\bm R_d)^\vee,\quad \e_\Omega=\bm\Omega-\bm R\T\bm R_d\bm\Omega_d .
\end{equation}
The wrench the cables must supply (\texttt{PlatformController::step}) is
\begin{align}
 \bm f&=m\Bigl(\bm a_d+k_p\e_p+k_d\e_v+k_i\!\!\int\!\e_p+g\e_3\Bigr)-\bm f_k-\hat{\bm d}_f,\\
 \bm\tau&=\bm R\Bigl[\bm J\Bigl(-k_p\e_R-k_d\e_\Omega-k_i\!\!\int\!\e_R-\bm\Omega\times\bm R\T\bm R_d\bm\Omega_d\Bigr)+\bm\Omega\times\bm J\bm\Omega\Bigr]-\bm\tau_k-\bm R\hat{\bm d}_\tau,
\end{align}
with $k_p=\omega_n^2$, $k_d=2\zeta\omega_n$, $k_i=0.5\,\omega_n^3$ ($\omega_n=7$ rad/s), bounded integrals, and the
observer's disturbance estimate fed forward. $\bm w=(\bm f,\bm\tau)$ goes through \eqref{eq:qp}.

\subsection{Hybrid winch commands}\label{sec:hybrid}
Drone winches run in tension mode \eqref{eq:tenmode} and ground winches in length mode \eqref{eq:lenmode}; all
get $T^{\mathrm{ff}}=\bm t$ from \eqref{eq:qp}.
\paragraph{Length-mode cables} servo to the inverse kinematics of a corrected pose,
$L^c_i=d_i(\bm p_d+k_{ik}\!\int\!\e_p,\ \bm R_d\exp\hatm{-k_{ik}\!\int\!\e_R},\ \hat{\bm a}_i)/(1+t_i/EA)$: in the
directions these cables hold stiffly a wrench integral has almost no authority, so the steady error is removed
kinematically.
They also carry a slow admittance on their load cells,
\begin{equation}
 L^c_i\leftarrow L^c_i+\delta L_i,\qquad \delta\dot L_i=k_a\,(T_i-t_i),\quad|\delta L_i|\le\delta_{\max},
 \label{eq:adm}
\end{equation}
so each is a stiff length servo above the bandwidth of \eqref{eq:adm} and follows its tension set-point below it.
This is what makes the scheme work with \emph{estimated} anchors: measured on the model, a millimetre of
ground-anchor error shifts a ground-cable tension by about 4 N, and \eqref{eq:adm} reduces the tension ripple from
2.2 to 0.7 N. With the true state it is switched off.
\paragraph{Tension-mode cables} get the rate of a commanded twist,
\begin{equation}
 \dot L^c_i=\bm u_i\T\bigl(\dot{\hat{\bm a}}_i-\bm v_c-\bm\omega_c\times\bm r_i\bigr),\qquad \bm v_c=\bm v_d+k_t\e_p,\quad\bm\omega_c=\bm R_d\bm\Omega_d-k_t\bm R\e_R,
 \label{eq:twist}
\end{equation}
with the correction limited to 0.05 m/s.
\paragraph{Why \eqref{eq:twist}.} Take one translational direction held by tension-mode cables with projection
$\gamma_i$, total damping $c=\sum_i\gamma_i^2c_w$. With $\dot L^c$ equal to the reference rate only
($k_t=0$), \eqref{eq:tenmode} and the PID give $c\,\dot e+mk_pe+mk_i\!\int\!e=0$ (the platform inertia is
negligible against $c$), i.e.
\begin{equation}
 \omega=\sqrt{mk_i/c},\qquad\zeta=\frac{mk_p}{2\sqrt{mk_ic}} .
\end{equation}
With $m=3$ kg, $k_i=171.5$, $k_p=49$, $c\approx1800$ N\,s/m: $0.085$ Hz and $\zeta=0.076$. The linearised
closed-loop map of the full model (Section~\ref{sec:stab}) has exactly this mode: 0.08 Hz, $\zeta=0.089$. With
\eqref{eq:twist} the same equation becomes $c(\dot e+k_te)+mk_pe+mk_i\!\int\!e=0$, with real poles near
$-k_t$ and $-mk_i/(ck_t+mk_p)$: the error decays at $k_t=3.5$ s$^{-1}$ and the oscillation is gone.

\subsection{Quadrotor}
Geometric tracking (Lee, Leok, McClamroch 2010) with the cable force fed forward from the load cell and the
estimated cable direction, $\hat{\bm F}_c=-T\hat{\bm u}-\tfrac12\rho gL\e_3$ (\texttt{DroneController}):
\begin{align}
 \bm F&=-\bm K_p(\bm p_q-\bm p_r)-\bm K_d(\bm v_q-\bm v_r)-\bm K_i\!\!\int\!(\bm p_q-\bm p_r)+m_q(\bm a_r+g\e_3)-\hat{\bm F}_c,\\
 f&=\bm F\!\cdot\!\bm R_q\e_3,\qquad \bm b_3=\bm F/\|\bm F\|,\quad\bm R_c=[\bm b_2\times\bm b_3,\ \bm b_2,\ \bm b_3],\quad \bm b_2=\frac{\bm b_3\times\bm b_{1c}}{\|\bm b_3\times\bm b_{1c}\|},\\
 \bm M&=-\bm K_R\,\tfrac12(\bm R_c\T\bm R_q-\bm R_q\T\bm R_c)^\vee-\bm K_\Omega\bm\Omega_q+\bm\Omega_q\times\bm J_q\bm\Omega_q-\bm r_a\times\bm R_q\T\hat{\bm F}_c ,
\end{align}
with the horizontal part of $\bm F$ limited to $F_z\tan\theta_{\max}$. The last term cancels the cable moment
about the centre of mass in \eqref{eq:quadrot}.

\subsection{Ground robot}
Desired planar velocity $\bm v^\star=\dot{\bm g}+k(\bm g-\bm x)$ towards the formation goal $\bm g$, limited to
0.6 m/s; heading error $e_\theta$ to $\bm v^\star$ (reversing if that turns less); $v=\|\bm v^\star\|\cos e_\theta$,
$\omega=k_\theta e_\theta$; acceleration-limited; wheel speeds $(v\mp\tfrac b2\omega)/r$ (\texttt{UgvController}).

\section{Perception and planning}\label{sec:plan}
\paragraph{Lidar and map} (\texttt{lidar\_scan}, \texttt{OccupancyGrid}). Each ground robot's 2-D lidar returns
ranges $r_k$ at bearings $\alpha_k$. A scan taken at the robot's \emph{estimated} pose updates a log-odds grid:
cells along the beam $\ell\leftarrow\ell-0.2$, the end cell $\ell\leftarrow\ell+0.9$, $|\ell|\le4$; a cell is
occupied for $\ell>0.7$. A two-pass chamfer transform gives the distance field $D(\bm x)$ to the nearest occupied
cell.
\paragraph{Team validity} (\texttt{TeamPlanner::valid}). With the formation centred on the platform at
$\bm c$ and ground robots at $\bm c+\bm o_k$, the position is valid iff
\begin{equation}
 D(\bm c)\ge r_b,\qquad D(\bm c+\bm o_k)\ge r_r,\qquad D(\bm c+s\bm o_k)\ge r_c\ \ \forall s\in(0,1),\ k=1..4,
\end{equation}
each radius enlarged by 8\% (chamfer error) plus one cell. The map is two-dimensional, so the team does not pass a
cable over an obstacle, and unmapped space counts as free.
\paragraph{Search.} A* over the grid (8-connected, octile heuristic) with lazily cached validity, then
line-of-sight shortcutting; it is repeated as the map grows.

\paragraph{Three dimensions} (\texttt{depth\_scan}, \texttt{ElevationGrid}, \texttt{TeamPlanner3}). A downward
depth camera on each drone gives points on whatever is below; the map keeps the highest return per cell,
$H(\bm x)$, and a lidar return raises its cell to at least the height of the scan plane. Dilating $H$ by the robot,
cable and tool radii gives $H_r,H_c,H_b$. A platform position $\bm c=(x,y,z)$ with the formation centred on it is
valid iff
\begin{align}
 &H_b(x,y)+\delta\le z-\ell_{\mathrm{tool}}, \qquad H_r\bigl((x,y)+\bm o_k\bigr)\le h_{\mathrm{step}},\\
 &H_c\bigl(\bm q_k(s)\bigr)+\delta\le z_k(s)\quad\forall s\in[0,1],\ \text{all eight cables},
\end{align}
where $\bm q_k(s),z_k(s)$ run along cable $k$ from its platform corner to its fairlead (0.35 m on a ground robot,
4.5 m on a drone) and $\delta$ is the vertical clearance. The search is A* over $(x,y,z)$ (26-connected) with the
metric $\|\Delta xy\|+c_z|\Delta z|$, so the team lifts the payload over a low obstacle while the ground robots
straddle it, and goes around a tall one. Cells within 0.3 m of the start are passable, so a start that new
measurements have just made invalid can be left.

\paragraph{Reconfiguration} (\texttt{TeamPlanner3::plan4}). The ground robots' ring may be scaled by
$\sigma\in\{0.6,\dots,1\}$, so the state is $(x,y,z,\sigma)$ and the robots sit at $(x,y)+\sigma\bm o_k$. A
$(z,\sigma)$ pair is admitted only if the tension distribution of Section~\ref{sec:td} meets the weight exactly
with every tension above $t_{\min}+0.5$ N (tabulated once: level platform, formation centred). Moves are the 26
spatial neighbours at constant $\sigma$ and a change of $\sigma$ in place, costing
$c_\sigma|\Delta\sigma|$. While the plan is followed the ring radius and the height change continuously, and the
winches follow with the cable lengths from the inverse kinematics.

\section{Keeping the cables taut}\label{sec:taut}
Three things were needed.
\paragraph{The hold region must be wrench-feasible.} The formation holds position while the platform is inside a
radius $R(z)$ and that radius was computed with a test that checked $\bm t_{\min}\le\bm t\le\bm t_{\max}$ but not
$\bm W\bm t=\bm w$. With cables pinned at a bound the others cannot always balance the wrench, so the platform was
allowed to go where a ground cable has to fall below $t_{\min}$. With the residual checked
($\|\bm W\bm t-\bm w\|<0.05$ N), the feasible radius at $z=1.3$ m is 0.40 m, not 0.88 m.
\paragraph{Reference governor} (\texttt{PlatformController::governed}). The commanded reference is followed by
$\dot{\bm v}_g=\mathrm{sat}_{a_{\max}}\bigl(\bm v_d+k_g(\bm p_d-\bm p_g)-\bm v_g\bigr)/\Delta t$, $\dot{\bm p}_g=\bm v_g$ (and the
same on $\SO$), so a velocity step in the command is not a tension impulse.
\paragraph{Slack guard.} A length-mode cable measured below $t_g$ is reeled in at $k_g^{T}(t_g-T_i)$ on top of
\eqref{eq:adm}.

\section{Stability}\label{sec:stab}
\subsection{Platform loop with ideal wrench}
If the cables deliver $\bm w$ exactly and $\hat{\bm d}=\bm d$, the translational error obeys
$\ddot{\e}_p+k_d\dot{\e}_p+k_p\e_p+k_i\!\int\!\e_p=0$, Hurwitz iff $k_dk_p>k_i$ (here $12.6\cdot49\gg171.5$). For the
attitude, with $\Psi=\tfrac12\mathrm{tr}(\bm I-\bm R_d\T\bm R)$ and
$V=\tfrac12\e_\Omega\T\bm J\e_\Omega+k_p\lambda_J\Psi+c_1\e_R\T\bm J\e_\Omega$, the argument of Lee et al.\ gives
$\dot V\le-\bm z\T\bm Q\bm z$, $\bm z=(\|\e_R\|,\|\e_\Omega\|)$, $\bm Q\succ0$ for small $c_1$, on the sublevel set
$\Psi<2$: exponential convergence from any initial attitude error below $180^\circ$.
\subsection{What the ideal analysis leaves out}
The wrench is produced through \eqref{eq:qp}, the winch servos \eqref{eq:lenmode}--\eqref{eq:tenmode}, elastic
cables and robots that are themselves feedback systems, at three sample rates. The time scales are separated
(cable axial mode 23 Hz, winch length servo 8 Hz, drone attitude about 3 Hz, platform loop 1.1 Hz, drone position
0.4 Hz) but not by enough for a clean singular-perturbation proof, and the hybrid winch scheme adds the slow mode
analysed in Section~\ref{sec:hybrid}.
\subsection{Linearised closed-loop map}
The sampled closed loop is therefore checked on the full model. Let $\Phi$ be the map of the closed-loop state
(plant states, controller integrators) over one common control period $T_c=20$ ms (20 physics steps; all three
controllers fire at its start). \texttt{Simulator::monodromy} computes $\bm A=\partial\Phi/\partial\bm x$ at an
equilibrium by central differences in error-state coordinates ($\delta\bm\theta$ for rotations). The
equilibrium is locally exponentially stable iff the spectral radius $\rho(\bm A)<1$; an eigenvalue $\mu$
corresponds to the continuous pole $s=\ln\mu/T_c$. The ground robots are parked for this test: a
differential-drive base cannot be asymptotically stabilised to a point by smooth feedback (Brockett), and its
go-to-goal law has a dead band, so its position is a neutral direction by design.

\section{Results}\label{sec:results}
\subsection{Model against the Python reference (\texttt{tests/core\_test.py}, part 1)}
Structure matrix: maximum difference $3\times10^{-16}$ over 300 random poses. Tension distribution: identical to
the closed form to $2\times10^{-12}$ N when no bound is active; projected gradient of the bounded solutions below
$10^{-13}$. Drone caps equal to the Python implementation. Forward kinematics recovers a pose to $10^{-9}$ m.

\subsection{Closed loop on the C++ plant (part 2)}
Lift 0.3 m with $10^\circ$ roll and $15^\circ$ yaw, then a 22 N payload.
\begin{center}\resizebox{\ifdim\width>\textwidth\textwidth\else\width\fi}{!}{\begin{tabular}{lcc}
\toprule
 & true state & sensors and estimators in the loop \\
\midrule
position error held after the move (mean / max) & 0.5 / 0.8 mm & 0.8 / 1.1 mm \\
attitude error (max) & 0.10$^\circ$ & 0.26$^\circ$ \\
under the payload (mean / peak) & 0.3 / 2.7 mm (announced) & 1.7 / 8.6 mm (not announced) \\
cable tensions & 6.7--35.5 N & 5.2--36.9 N \\
platform estimate error (mean / max) & -- & 1.3 / 4.1 mm, 0.78$^\circ$ max \\
payload estimated by the observer & -- & $-22.0$ N (true $-22$) \\
\bottomrule
\end{tabular}}\end{center}
The full loop (plant, sensors, estimators, controllers at 1 kHz) takes 10--13 $\mu$s per step, about 80 times
real time on one core.

\subsection{Stability (part 3)}
Linearised map over 20 ms at the level pose, all loops closed, ground robots parked, 122 states.
\begin{center}\resizebox{\ifdim\width>\textwidth\textwidth\else\width\fi}{!}{\begin{tabular}{lccc}
\toprule
 & $\rho(\bm A)$ & slowest pole & least-damped mode \\
\midrule
original hybrid scheme ($k_t=k_{ik}=k_a=0$) & 0.99990 & $-0.005$ s$^{-1}$ & 0.08 Hz, $\zeta=0.089$ \\
with commanded twist \eqref{eq:twist} and pose-integral inverse kinematics & 0.99823 & $-0.089$ s$^{-1}$ & 22.8 Hz, $\zeta=0.15$ \\
as deployed (plus admittance \eqref{eq:adm}) & 0.99944 & $-0.028$ s$^{-1}$ & 22.9 Hz, $\zeta=0.15$ \\
\bottomrule
\end{tabular}}\end{center}
All eigenvalues are inside the unit circle. The $-0.089$ s$^{-1}$ pole is the drones' position integrator
($k_i/k_p$); the $-0.028$ s$^{-1}$ pole is the admittance redistributing the ground-cable lengths; the 23 Hz mode
is the cables' axial mode. This is a local result at one equilibrium, for the model, with the true state fed
back; it is not a proof for the estimator-in-the-loop system or for large motions.

\subsection{In Isaac Sim}
\texttt{tests/isaac\_physics\_test.py}: the same manoeuvre, then a 1 m translation of the whole formation;
statistics over every physics step after 1 s. The position error is measured against the commanded reference,
so for the C++ core it includes the lag of the reference governor at the two command steps (the maxima).
\begin{center}\resizebox{\ifdim\width>\textwidth\textwidth\else\width\fi}{!}{\begin{tabular}{lccc}
\toprule
 & Python controllers, & C++ core, & C++ core, \\
 & true state & true state & sensors in the loop \\
\midrule
position error rms / max & 4.0 / 11.4 mm & 2.1 / 13.2 mm & 3.0 / 14.2 mm \\
attitude error mean & 0.71$^\circ$ & 0.26$^\circ$ & 0.77$^\circ$ \\
a cable slack & 0.55\% of the time & never & never \\
lowest tension & 0 N & 1.4 N & 3.0 N \\
platform estimate error mean / max & -- & -- & 1.3 / 2.8 mm \\
\bottomrule
\end{tabular}}\end{center}
\texttt{run\_traj.py --scenario track --core} (figure-eight, 48 s, sensors in the loop): 4.4 mm mean, 15.9 mm max;
a cable slack 0.15\% of the time (about 70 ms in total).

\subsection{Sensor-based planning}
\texttt{tests/core\_plan\_test.py} (2-D, lidar only, no prior map, replanning every 0.5 m): goal reached in 34
plans of about 3 ms; path 16.5 m, around both obstacles; true clearances 0.55 m for the ground robots (0.55
required) and for the cables (0.20).

\texttt{tests/core\_plan3d\_test.py} (elevation map from depth cameras and lidars): goal reached in 31 plans of
about 21 ms; path 15.4 m, around the cabin and over the pallet at $z=1.5$ m with the ground robots straddling it;
true clearances 0.81 m for the ground robots, 0.32 m for the cables, 0.22 m for the tool.

\texttt{run\_traj.py --scenario plan3d --core} (the same planner running in Isaac while the team moves, sensors
and estimators in the loop): goal reached, 15.4 m in 57 s, 31 plans of 21 ms; tracking error 3.1 mm mean, 25 mm
max at the corners of the plan; recorded clearances 0.80 m for the ground robots and 0.34 m for the cables; a
cable slack 0.04\% of the time. The first attempt drove a ground robot into the stack of blocks on the site,
which was not among the obstacles given to the range models; it is now.

\texttt{run\_traj.py --scenario course --core} (reconfiguration; a 4.4 m gate in a 2 m wall for a team 5.0 m wide,
a crate, two 3 m pillars; search over $(x,y,z,\text{ring scale})$ with the wrench-feasible (height, scale) pairs
tabulated beforehand): goal reached, 19.4 m in 68 s, 39 plans of 92 ms. The ground robots' ring goes from 2.75 m to
1.65 m at the gate and back, and the ground cables from 2.96 m to 1.74 m (drone cables 3.49--3.62 m). Recorded
clearances: ground robots 0.80 m, cables 0.80 m, drones 1.73 m. Tracking error 2.8 mm mean, 25 mm max; no cable
slack (lowest tension 0.4 N).

\subsection{Limits}
The sensors are measurement models on the true simulator state: ranges are cast against the collider boxes, tag
poses and fixes are the true values with noise; no images or point clouds are processed and there is no latency.
The sensors do not see the team itself. The platform has no inertial sensor, so its attitude estimate
(0.3--0.9$^\circ$) is the weakest part. The planner assumes the formation centred on the platform and a level
platform, and treats unseen space as free. The cables are never slack in the manoeuvre test but still are for a
few tens of milliseconds in the two long trajectory runs.


\end{document}
